\section{Contexte}\label{sec:contexte}
Le principe de l’alignement apparaît déjà sur deux courtes séquences. Dans \texttt{AC} et \texttt{AGC}, les deux \texttt{C} occupent des positions différentes. Pour les mettre en regard, il suffit de laisser dans la première séquence un espace à la hauteur du \texttt{G} de la seconde. Cet espace représente la seule insertion nécessaire pour passer de \texttt{AC} à \texttt{AGC}. Les deux articles étudiés donnent le même coût à une insertion, à une suppression ou à une substitution. Ils diffèrent par la manière de rechercher la suite d’opérations la moins coûteuse.

Sadiq et ses coauteurs (2020) s’appuient sur la programmation dynamique. Leur méthode compare les débuts des deux séquences, appelés \emph{préfixes}, en augmentant progressivement leur longueur~\cite[alg.~1 et fig.~4]{sadiq}. Dans la figure~\ref{fig:contexte-nvpd}, chaque case donne la distance entre le préfixe indiqué à gauche et celui qui figure en haut. Dans le calcul classique, la valeur inscrite à gauche doit être connue avant de passer à la colonne suivante. NvPD évite cette attente en repérant à l’avance la dernière occurrence de chaque lettre dans les préfixes de la seconde séquence. Si la lettre recherchée est présente, on connaît ainsi le nombre de lettres qu’il resterait à insérer après la correspondance. Le score de départ se trouve sur la ligne précédente, comme pour une substitution ou une suppression. Les colonnes de la nouvelle ligne peuvent donc être calculées en parallèle.

\schema{contexte-nvpd}{La flèche indique le calcul de la distance entre \texttt{A} et \texttt{AGC}. Les deux \texttt{A} se correspondent. Depuis le préfixe vide \(\varepsilon\), le coût est donc celui des deux insertions, soit \(D_{1,3}=D_{0,0}+2=2\).}{fig:contexte-nvpd}{report/figures/contexte-nvpd.tex}

Aguado-Puig et ses coauteurs (2022) prennent pour point de départ le nombre d’éditions que l’on autorise~\cite[alg.~1 à~4 et fig.~1 à~3]{wfagpu}. Tant que les lettres se correspondent, on peut avancer dans les deux séquences sans augmenter le coût de l’alignement. Chaque diagonale de la grille correspond à un décalage entre les deux séquences. Pour un coût donné, l’algorithme retient la position la plus avancée sur chaque diagonale accessible. Ces positions forment un \emph{front} (figure~\ref{fig:contexte-ewfa}). Le premier front qui atteint la fin des deux séquences donne leur distance. Sur GPU, chaque bloc peut traiter une paire et répartir ses diagonales entre ses fils de calcul.

\schema{contexte-ewfa}{Les deux \texttt{A} se correspondent au coût zéro. L’insertion de \texttt{G} ouvre la diagonale \(k=-1\), où la correspondance des deux \texttt{C} conduit au point \((2,3)\), à la fin des deux séquences.}{fig:contexte-ewfa}{report/figures/contexte-ewfa.tex}

Les essais publiés éclairent deux usages de l’alignement. Sur une paire de séquences de 1,02 et 1,17~Mnt, Sadiq et coll. mesurent le temps nécessaire pour obtenir la distance seule. Leur GPU GTX TITAN Black effectue ce calcul 92,04 fois plus vite que leur référence séquentielle sur Core~i7~\cite[tabl.~3]{sadiq}. Aguado-Puig et coll. mesurent aussi le temps de reconstruction des alignements. Pour dix millions de paires de 1\,000 bases à 5\,\% d’erreurs, les chemins sont obtenus en 689~ms sur V100, transferts compris, contre 3\,900~ms avec eWFA sur les 80 fils du POWER9~\cite[tabl.~1]{wfagpu}.

\synthese{Dans cet exemple, NvPD remplit les six cases intérieures de la grille. eWFA-GPU atteint la dernière case dès la première édition autorisée. Le protocole examine comment cet écart de travail évolue avec la longueur et la divergence des séquences.}
